\section*{September 7, 2026: School definitions and raw housing trends}
\textit{Codex draft; data review and descriptive presentation requested by Jacob.}

The data review reproduces the existing 10,921-sale comparison: 3,287
transactions near selected closed schools and 7,634 near controls. The new
descriptive packet begins with the school roster and sample construction,
then shows annual prices, sales counts, property composition, and neighborhood
coverage. Each transaction receives equal weight. The school selection and
price rules remain unchanged. This entry records the initial packet with
complete characteristics; the later entry expands the descriptive sample.

All 129 names in the supplied roster match the original CPS February list,
recovered from a contemporaneous news archive. Appendix A of the Consortium's
2015 \textit{School Closings in Chicago} matches all 47 closed programs and
their 53 receiving-school assignments. Fifteen closed programs had receiving
schools move into their buildings. CPS notices confirm that Fermi/South Shore
and Garfield Park/Faraday already shared buildings. Excluding these 17 cases
produces the 30 closed programs used for housing, at 29 physical sites.
This selection reflects continued school use after the closure decision;
it does not establish subsequent vacancy or measure all program closures.

The 49 controls also require explanation. Of the 82 candidates that did not
close in 2013, the roster excludes 18 receiving schools, five turnarounds,
two delayed closures, four canceled closure proposals, three co-location
cases, and Mason, whose high-school program ended. All 49 control IDs appear
in the 2013--14 report card. The original notes govern the other exclusions;
individual final board reports for every action have not all been reverified.
In particular, whether to exclude the four canceled proposals remains a
substantive comparison-group question for discussion, not a cleaning error.

The housing checks independently reproduce every citywide price-sample ID,
the monthly CPI adjustment, and distances and nearest-school assignments for
all 239,254 study-period transactions. The separate finalization audit also
rechecks corrected property fields against its independent exemption
reconstruction. The price sample retains 87.7\% of market-screened sales near
closed schools and 90.8\% near controls. These losses combine property-type,
building-record, completeness, and price restrictions. Sales counts are reported
before those restrictions as well; without housing-stock denominators, they
are counts rather than turnover rates.

{\small Verification: All 103 existing script/build files were read, including
archived analyses. Runtime checks use the recorded housing sources; the
finalization check retains its existing post-exemption follow-up extract.
The school cleaner now rejects duplicate report-card IDs, and a new report
describes its unchanged 129-row output. Exemption downloads now publish only
completed transfers; disposable failure checks confirm that an interrupted
transfer preserves any existing snapshot. A stale housing-report statement
about school distances was corrected. Source hashes and the full school
exclusion lists are in the data review.}

\clearpage
\begin{center}
\includegraphics[width=\textwidth]{../input/overview_complete_prices.png}
\end{center}

From 2012 to 2018, mean transaction prices near closed schools rise from about
\$239,000 to \$307,000 in 2022 dollars, while medians rise from \$164,000 to
\$181,000. Near controls, means rise from \$149,000 to \$234,000 and medians
from \$61,000 to \$143,000. The different movements of means and medians are
already visible in the raw data. These changing sets of transactions combine
price changes with changes in the properties and neighborhoods represented;
the plotted differences do not establish a closure effect.

{\small Reproduction: \texttt{make -C tasks/audits/home\_data\_overview/code}
builds the ten-page \texttt{data\_overview.pdf}, individual figures,
\texttt{data\_review.md}, and saved-summary reports. The packet includes price
percentiles and within-property-type means alongside this figure. The review
uses the September 7 working tree and the source hashes recorded in
\texttt{data\_review.rds}; housing vintages and analytical choices are unchanged.}
